\documentclass[11pt]{article}
\usepackage{times}
\usepackage{fullpage}
\usepackage{url}
\usepackage{hyperref}
\usepackage{fancyhdr}
\usepackage{graphicx}
\usepackage{enumitem}
\usepackage{subcaption}
\usepackage{amsmath, amsfonts, amsthm, fouriernc}
\usepackage{IEEEtrantools}
\usepackage{parskip}
\usepackage{booktabs}

\pagestyle{fancy}
\addtolength{\headheight}{2em}
\addtolength{\headsep}{1.5em}
\lhead{\large{ME 2801}}

\usepackage{sectsty}
\sectionfont{\fontsize{12}{8}\selectfont}

\begin{document}

\begin{center}
  \Large{\bf{Assignment: Steady-State Error}}
\end{center}

This assignment exercises the steady-state error analysis from the chapter.
All computations assume \emph{unity negative feedback} and a \emph{stable}
closed loop unless stated otherwise.  MATLAB is encouraged throughout as a
verification tool --- a sample script is in
\texttt{code/assignment\_solutions.m}.

% ---------------------------------------------------------------
\section*{Problem 1 --- Three test inputs}
% CLAUDE: Change "Problem" to "Exercise" for all.
The unity-feedback system has open-loop transfer function
\[
  G(s) = \frac{600}{(s+10)(s^2 + 4s + 25)}.
\]

% CLAUDE: Is there any type of engineering system that a transfer function with this form of dynamics: 1 real pole, 2 complex and no zeros?  If so, give an example in one sentence.  If not, explain why not.  Maybe this is both a plant and controller for something?   Good to highlight G(s) in this context is both controller and plant.


\begin{enumerate}[label=(\alph*)]
  \item Identify the system Type by inspection of $G(s)$.
  \item Compute the static error constants $K_x$, $K_v$, and $K_a$. % CLAUDE: I don't like these as ane end.  We treat these as just named variables - intermediate steps in the general expression and used in the process of finding the steady-state error
  \item Find the steady-state error for the inputs $u(t)$, $t\,u(t)$, and
    $\tfrac{1}{2}\,t^2 u(t)$. % CLAUDE: Provide a little note about why the factor of 1/2 is there.  It's just to make the math simpler, but it might be confusing to students.
  \item In MATLAB, verify your step-input answer in (c) two ways:
  \begin{itemize}
    \item Calculate: compute the closed-loop transfer function and evaluate its DC gain.
    \item Simulate: generate the step response and check the final value of the error signal.
  \end{itemize}
  \item Two of your answers in (c) are infinite.  In one sentence, explain
    what physical fact about $G(s)$ makes that so.
\end{enumerate}

% ---------------------------------------------------------------
\section*{Speed (transient) and accuracy (steady-state) together}

The unity-feedback system has an open-loop transfer function
\[
  G(s) = \frac{4000}{s(s+60)}.
\]

\begin{enumerate}[label=(\alph*)]
  \item What is the percent overshoot for a unit step input?  (Recall the
    second-order formulas from the time-response chapter.) % CLAUDE: Open-loop or closed-loop?  be specific
  \item What is the 2\,\% settling time?
  \item Find the steady-state error for the inputs $5u(t)$, $5tu(t)$, and
    $5t^2 u(t)$. % CLAUDE:  somehow say they need to do this analytically, by-and 
  \item Verify (a)--(c) in MATLAB using \texttt{stepinfo} and a closed-loop
    \texttt{lsim} on the ramp. 
    % CLAUDE: Refactor the above where the question is to dothe same thing twice, once with pen and paper and once with matlab.
  \item Suppose you want to reduce the ramp error in (c) by raising the loop
    gain.  In one or two sentences, explain why this might be a bad idea ---
    referencing your answer to (a).
\end{enumerate}

% ---------------------------------------------------------------
\section*{Problem 3 --- Design a gain to meet a steady-state spec}

The forward-path transfer function of a unity-feedback position-control
system is
\[
  G(s) = \frac{K(s+5)}{s\,(s+8)(s+12)}.
\]

% CLAUDE: Phrase this as designing a proportional only controller.  Give them to TFs, plant (G(s)) and controller (C(s) = K_p).  Add a hint to keep track of what is the open-loop (forward path) transfer function and what is the closed-loop transfer function.  

\begin{enumerate}[label=(\alph*)]
  \item What is the system Type?
  \item Find the value of $K$ such that the steady-state error to a unit
    ramp input is 4\,\%.
  \item Use MATLAB to compute the closed-loop poles for your value of $K$.
    Are they all in the left half-plane?
  \item What single assumption from the chapter makes (b) a meaningful
    answer?  In one sentence, what would you check before deploying this
    controller?
\end{enumerate}

% ---------------------------------------------------------------
\section*{Problem 4 --- Fill the table}
% CLAUDE: comment out this section.
For each of the two open-loop transfer functions below in unity feedback, fill in the steady-state
error: a number, $0$, or $\infty$.
\[
  G_a(s) = \frac{20}{(s+2)(s+10)},
  \qquad
  G_b(s) = \frac{40}{s\,(s+8)}.
\]

\begin{enumerate}[label=(\alph*)]
  \item Identify the Type of each plant before computing.
  \item Fill in the six cells.
  \begin{center}
  \renewcommand{\arraystretch}{1.5}
  \begin{tabular}{l|ccc}
    \toprule
    & $r(t) = u(t)$ & $r(t) = t\,u(t)$ & $r(t) = \tfrac12 t^2 u(t)$ \\
    \midrule
    $G_a(s)$ & \rule{2.5cm}{0pt} & \rule{2.5cm}{0pt} & \rule{2.5cm}{0pt} \\
    $G_b(s)$ & & & \\
    \bottomrule
  \end{tabular}
  \end{center}
  \item In one sentence, what is the structural difference between $G_a$ and
    $G_b$ that explains the entire pattern in your table?
\end{enumerate}

% ---------------------------------------------------------------
\section*{Problem 5 --- USV Surge: Predict before you simulate} % CLAUDE: Make all section titles use this capitalization style.

The USV surge (forward-speed) plant, identified from the lab data, is
\[
  G_\mathrm{surge}(s) = \frac{5.1}{1.3\,s + 1}.
\]
For each controller architecture below, predict the steady-state error
analytically for a unit-step (1\,m/s) speed command.  Then verify each
prediction in MATLAB by computing the closed-loop transfer function and
plotting the step response.

\begin{enumerate}[label=(\alph*)]
  \item \textbf{P only}: $K_P = 0.5$.
  \item \textbf{PI}: $K_P = 0.5$, $K_I = 0.4$.
  \item \textbf{P + feedforward}: $K_P = 0.5$, $K_{\!f\!f} = 1/K_\mathrm{plant}$.
  \item \textbf{PI + feedforward}: same $K_P, K_I$ as (b), and $K_{\!f\!f}$ as (c).
  \item Use the \emph{adaptation argument} from the chapter (not the
    steady-state error formula) to explain in one or two sentences why the
    PI controller in (b) drives the error to zero.
\end{enumerate}

% ---------------------------------------------------------------
\section*{Problem 6 --- USV yaw: ramp tracking design}

The USV's yaw-rate plant, identified from the lab data, is
\[
  G_\mathrm{yaw\text{-}rate}(s) = \frac{0.82}{0.61\,s + 1}.
\]
Heading is the integral of yaw rate, so the plant from rudder input to
\emph{heading angle} picks up an extra factor of $1/s$:
\[
  G_r(s) = \frac{0.82}{s\,(0.61\,s + 1)}. % CLAUDE: Have students find this.
\]
A proportional controller $C(s) = K_r$ is used in a unity-feedback heading
loop.  The heading reference ramps at $5^\circ$/s (a constant turn-rate
command).

\begin{enumerate}[label=(\alph*)]
  \item What is the system Type of the closed loop?  Justify in one phrase.
  \item Find the velocity error constant $K_v$ in terms of $K_r$, and write
    the steady-state heading lag in degrees as a function of $K_r$.
  \item The design spec is a heading lag $\le 1^\circ$.  Find the minimum
    $K_r$ that meets the spec.
  \item Verify your $K_r$ in MATLAB by simulating the closed-loop ramp
    response and checking the steady-state error.
  \item Your collaborator says: ``The plant already has a $1/s$ in it, so
    it's Type~1.  Increasing $K_r$ makes it Type~2, which gives us zero
    ramp error.''  Critique this statement in one or two sentences.
\end{enumerate}

\end{document}
